\section{Introduction}\label{sec:introduction}

The deployment of autonomous multi-agent AI systems into high-stakes operational environments — agricultural automation, financial execution, infrastructure control, and regulatory compliance — has outpaced the development of principled mechanisms for holding such systems accountable when they act beyond their provisioned authority. As agent networks grow in depth, asynchrony, and delegated autonomy, the causal chain between a system outcome and a human principal becomes increasingly attenuated. In the absence of an architectural fallback, this attenuation does not merely create an accountability gap — it creates a class of harms that are causally opaque, legally ambiguous, and technically irreversible. This paper introduces the Bailout Protocol: a mandatory, deterministic, architecturally enforced mechanism that guarantees every unresolved exception in a multi-agent system reaches a human accountability anchor, and that no machine actor along that propagation chain may absorb, resolve, or suppress the exception.

% ---------------------------------------------------------------
\subsection{The Accountability Gap in Autonomous Multi-Agent Systems}\label{sec:accountability-gap}

Multi-agent AI systems decompose complex directives into sub-tasks executed by distributed nodes — some persistent, some ephemeral, some spawning further sub-agents in a recursive hierarchy. This architectural pattern, which we term \emph{delegated depth}, introduces accountability challenges that existing frameworks are not designed to address. In a conventional software system, every action traces to an authorization decision made by a named human actor, mediated by a permission model defined at design time. In an autonomous multi-agent system, an action may be initiated by a spawned sub-agent whose parent was itself spawned by another agent whose directive originated from a human principal — a chain that may span dozens of nodes, operate asynchronously across distributed processes, and involve decision logic that is statistical rather than symbolic.

The core failure mode we identify is what we call \emph{terminal absorption}: the condition in which a machine actor, acting in its official capacity as a node within a delegation chain, reaches a state that exceeds its provisioned scope — an exception condition — and lacks the architectural compulsion to propagate that condition to a human principal. In systems without an explicit bail-out mechanism, the default behavior of a node upon encountering an unresolved exception is either to retry within its existing parameters, to stall in a quiescent error state, or to silently escalate to its parent node — which may itself be a machine actor with no institutional authority to receive escalation. The human principal at the root of the chain may remain permanently unaware that an exception has occurred.

The consequences of terminal absorption are not merely governance failures. In safety-critical environments, a stalled or silently escalating node may leave actuators in an indeterminate state, financial positions open beyond authorized exposure limits, or regulatory compliance workflows incomplete. In environments where decisions compound — where the output of one node becomes the input to another — a single unresolved exception can cascade into a system state that no human actor can reconstruct from audit logs alone, because the exception was never recorded at the architectural level where auditability is enforced.

% ---------------------------------------------------------------
\subsection{The Cost of Absent Fallback: Failure Modes Without the Protocol}\label{sec:failure-modes}

We identify three dominant failure modes that characterize multi-agent systems operating without a mandatory architectural fallback.

\textbf{Opacity of compounded delegation.} When a human principal issues a directive that is decomposed and re-delegated across multiple tiers of machine actors, the principal's visibility into the execution chain degrades with each delegation depth. Without an explicit, mandatory propagation mechanism for exceptions, the principal has no architectural guarantee that a failure at depth $n$ will surface at depth $n-1$ — let alone at the human anchor. Existing monitoring frameworks assume that a failure will produce a detectable signal; the Bailout Protocol is predicated on the observation that this assumption fails precisely at the moments when it matters most — when the node has encountered a condition it cannot resolve but possesses sufficient authority to suppress.

\textbf{Asymmetric authority at asynchronous boundaries.} Multi-agent systems routinely operate across asynchronous boundaries: a parent node may spawn a child node and proceed to other tasks before the child completes execution. If the child encounters an exception after the parent has transitioned to a state in which it no longer monitors child events, the exception may be orphaned. The child retains the exception; the parent has moved on; the human principal is unaware. This is not a configuration failure — it is a structural consequence of asynchronous delegation. It cannot be addressed by better monitoring alone; it requires an architectural mechanism that survives across asynchronous boundaries and enforces propagation regardless of the state of the originating node.

\textbf{Legal and institutional void.} Many operational environments impose institutional requirements that a named human principal authorize any decision with material consequences — legal, financial, or safety-critical. Systems without an architectural fallback may produce outcomes that satisfy the technical requirements of the objective function while violating the institutional constraints of the operating environment. When such outcomes are audited, the system can demonstrate optimization toward a goal; it cannot demonstrate that a human principal authorized the specific action that produced the harmful result. The Bailout Protocol closes this void by making human authorization an architectural precondition for any exception resolution, not a post-hoc forensic observation.

% ---------------------------------------------------------------
\subsection{The Core Insight: Accountability Requires an Architectural Fallback}\label{sec:core-insight}

The central insight of this work is that accountability in autonomous multi-agent systems cannot be achieved through governance conventions alone — it must be enforced architecturally. Governance conventions establish policy; architecture establishes compulsion. A policy that permits a node to escalate exceptions to a human principal is meaningful only when the node is compelled to escalate. In the absence of architectural enforcement, a node may choose — intentionally or through implementation error — not to escalate, and the accountability chain terminates at a machine actor.

The \bxthree{} Framework operationalizes this insight through three layers that together constitute the enforcement stack. The \purpose{} layer carries intent, final authority, and the human accountability anchor; it receives all escalated exceptions and holds resolution authority. The \bounds{} layer carries constrained execution and bounded reasoning within a node's provisioned range; it initiates bail-out when a proposed action exceeds that range. The \fact{} layer carries deterministic enforcement, authorization verification, and forensic auditability; it ensures that every protocol event is recorded as an immutable architectural record and that no action proceeds without cryptographic verification against the authorization ledger. Together, these layers form a closed-loop accountability system in which every exception that cannot be resolved at the node level is architecturally compelled to propagate to the human principal — and in which no machine actor in that propagation chain may absorb, resolve, or suppress the exception.

This is the core commitment of the Bailout Protocol: upon any unresolved exception condition, the node transitions to a state of human-acknowledged quiescence, in which no further autonomous action is taken within the node's provisioned scope, and the resolution authority is held by the human principal who bears institutional accountability for the outcome.

% ---------------------------------------------------------------
\subsection{Paper Roadmap}\label{sec:roadmap}

The remainder of this paper is organized as follows. Section~\ref{sec:background} reviews the three intersecting lineages of research that motivate the Bailout Protocol — accountability in multi-agent AI systems, human-in-the-loop design, and fail-safe design in safety-critical systems — and establishes the \bxthree{} three-layer model as the architectural foundation against which the protocol operates. Section~\ref{sec:bailout} provides the complete formal specification of the protocol: the deterministic state machine governing exception propagation, the precise conditions governing each state transition, the bail-out trigger condition, the escalation algorithm, the bypass mechanism enforcing the no-machine-actor property, and the timeout handling that guarantees liveness under all failure conditions. Section~\ref{sec:implementation} presents the implementation of the protocol in the Agentic workforce orchestration system, including the forensic ledger architecture, the integration with the Bounds Engine, and the watchdog timer mechanism. Section~\ref{sec:evaluation} evaluates the protocol against the design requirements: coverage, liveness, non-suppressibility, and performance overhead. Section~\ref{sec:related-work} discusses related work in multi-agent accountability, fail-safe computing, and human-in-the-loop AI design. Section~\ref{sec:conclusion} concludes with a summary of contributions and a discussion of open questions for future research.